% ==========================================================================
%  Chapter 80 -- Embedded rebar with penalty coupling and fibre-truss
%                feedback in sub-models FE²
% ==========================================================================

\chapter{Acero embebido con acoplamiento por penalidad e interacción
         fibra-truss en sub-modelos FE\textsuperscript{2}}
\label{ch:embedded-rebar}

% ------------------------------------------------------------------
\section{Motivación}
\label{sec:emb-motivation}

En el esquema multiescala FE\textsuperscript{2} descrito en capítulos
anteriores (\ref{ch:two-way-fe2}, \ref{ch:multiscale-16storey}),
las barras de refuerzo dentro de los sub-modelos continuos se
discretizaban como elementos truss (línea) cuyos nodos coincidían con
los nodos de esquina de la malla hexaédrica.  Esto producía dos
problemas:

\begin{enumerate}
  \item \textbf{Posición no realista:}
        Para una malla gruesa ($N_x = N_y = 2$), la función
        \texttt{std::round} mapeaba la coordenada física de la barra
        al índice de nodo entero más cercano.  Una barra diseñada a
        $y_0 = -0{.}200$\,m (\textit{cover} + $\phi/2$) con
        $\Delta x = 0{.}250$\,m se redondeaba a $i_x = 0$, es decir,
        sobre la cara \emph{exterior} del elemento.  Esto es
        físicamente incorrecto porque ignora el recubrimiento de
        concreto y puede representar una condición de contorno
        diferente a la real.

  \item \textbf{Acoplamiento artificial:}
        Al compartir nodos con el hexaedro, la barra heredaba la
        cinemática completa del concreto (3 grados de libertad
        compartidos), lo que impedía capturar el deslizamiento relativo
        barra--concreto y la contribución rigidizante correcta del
        acero.
\end{enumerate}

La solución implementada consiste en:
\begin{itemize}
  \item Crear \emph{nodos independientes} para las barras en sus
        posiciones físicas exactas dentro de la sección transversal.
  \item Acoplar los nodos de la barra al continuo hexaédrico mediante
        \emph{resortes de penalidad} que reproducen la interpolación
        Master-Slave.
  \item Incorporar el diámetro real de la barra en la exportación VTK
        para visualizar tubos con espesor físico en ParaView.
\end{itemize}


% ------------------------------------------------------------------
\section{Posición física de las barras}
\label{sec:emb-position}

Cada barra se especifica con sus coordenadas locales $(l_y, l_z)$
relativas al centroide de la sección, su área $A_b$ y su diámetro
$\phi_b$.  La estructura \texttt{RebarBar} se define como:

\begin{verbatim}
struct RebarBar {
    double ly{0};          // coordenada local y [m]
    double lz{0};          // coordenada local z [m]
    double area{0};        // área [m²]
    double diameter{0};    // diámetro [m]
    std::string group = "Rebar";
};
\end{verbatim}

Para una columna cuadrada de $b = h = 0{.}50$\,m con recubrimiento
$c = 40$\,mm y barras de $\phi = 20$\,mm, las coordenadas extremas son:
%
\begin{equation}
  y_0 = -\frac{b}{2} + c + \frac{\phi}{2}
      = -0{.}250 + 0{.}040 + 0{.}010 = -0{.}200\;\text{m}
\end{equation}
%
lo que sitúa la barra a 50\,mm del borde exterior, dentro del acero
real de la sección.  Las ocho barras (4 esquina + 4 central de cara)
se posicionan según:
%
\begin{center}
\begin{tabular}{ccc}
\toprule
Barra & $l_y$ [m] & $l_z$ [m] \\
\midrule
1--4 (esquinas) & $\pm 0{.}200$ & $\pm 0{.}200$ \\
5--6 (caras $z$) & $0{.}000$ & $\pm 0{.}200$ \\
7--8 (caras $y$) & $\pm 0{.}200$ & $0{.}000$ \\
\bottomrule
\end{tabular}
\end{center}


% ------------------------------------------------------------------
\section{Nodos independientes de barra y datos de imbricación}
\label{sec:emb-nodes}

La función \texttt{make\_reinforced\_prismatic\_domain()} genera la
malla hexaédrica (idéntica a \texttt{make\_prismatic\_domain()}) y
luego crea nodos de barra adicionales a lo largo del eje $z$.  Cada
barra produce $n_z \cdot s + 1$ nodos ($s = 1$ para Hex8, $s = 2$ para
Hex27), con identificadores que comienzan después de los nodos de la
malla hex:
%
\begin{equation}
  \text{nid}_{b,k} = N_{\text{hex}} + b \cdot
    (s \cdot n_z + 1) + k, \qquad k = 0, \ldots, s \cdot n_z,
\end{equation}
%
donde $N_{\text{hex}}$ es el número total de nodos hexaédricos.

Para el caso de la simulación del edificio ($n_z = 8$, Hex27,
8~barras), se crean $8 \times 17 = 136$ nodos de barra y $8 \times 8
= 64$ elementos línea (truss).  El total de nodos del sub-modelo pasa
de 850 (sólo hex) a 986.

\subsection{Elemento huésped y coordenadas padres}

Para cada nodo de barra, se identifica el hexaedro que lo contiene
(\emph{host element}) y se calculan las coordenadas padres
$(\xi, \eta, \zeta) \in [-1,1]^3$ mediante:
%
\begin{align}
  i_x^{\text{host}} &= \text{floor}\!\!\bigl((l_y + w/2) / \Delta x\bigr),
  &\xi  &= 2\,\frac{(l_y + w/2)/\Delta x - i_x^{\text{host}}}{1} - 1, \\[4pt]
  i_y^{\text{host}} &= \text{floor}\!\!\bigl((l_z + h/2) / \Delta y\bigr),
  &\eta &= 2\,\frac{(l_z + h/2)/\Delta y - i_y^{\text{host}}}{1} - 1, \\[4pt]
  i_z^{\text{host}} &= \text{floor}(k_z / s),
  &\zeta &= 2\,\frac{k_z - i_z^{\text{host}} \cdot s}{s} - 1,
\end{align}
%
donde $k_z$ es el sub-nivel $z$ del nodo de barra ($0 \le k_z \le s
\cdot n_z$).  Esta información se almacena en la estructura
\texttt{RebarNodeEmbedding}:

\begin{verbatim}
struct RebarNodeEmbedding {
    PetscInt rebar_node_id;
    int host_elem_ix, host_elem_iy, host_elem_iz;
    double xi, eta, zeta;
};
\end{verbatim}


% ------------------------------------------------------------------
\section{Condiciones de contorno en los extremos de las barras}
\label{sec:emb-bc}

Los nodos de barra que caen en las caras $z_{\min}$ y $z_{\max}$ del
sub-modelo reciben condiciones de Dirichlet computadas a partir de la
cinemática seccional del elemento viga (Timoshenko) en las posiciones
$(l_y, l_z)$ exactas de la barra:
%
\begin{equation}
  \mathbf{u}_{\text{bar},\text{face}} =
    \mathbf{u}_c + \boldsymbol{\theta} \times (l_y\,\hat{e}_y +
      l_z\,\hat{e}_z),
\end{equation}
%
donde $\mathbf{u}_c$ y $\boldsymbol{\theta}$ son el desplazamiento y
la rotación del centroide de la sección, evaluados en el extremo A
($z_{\min}$) o B ($z_{\max}$).

Estos nodos se agregan a las listas \texttt{face\_min\_z\_ids} y
\texttt{face\_max\_z\_ids} en la función
\texttt{build\_sub\_models()} del \texttt{MultiscaleCoordinator},
inmediatamente después de \texttt{cache\_face\_nodes()}.


% ------------------------------------------------------------------
\section{Acoplamiento por penalidad (Master-Slave)}
\label{sec:emb-penalty}

Los nodos interiores de barra (todos excepto los de las caras $z$)
no están conectados a la malla hexaédrica por el sieve de PETSc.
Para acoplar su cinemática al continuo se emplea un esquema de
\emph{penalidad}, equivalente a Master-Slave con interpolación por
funciones de forma.

\subsection{Formulación}

Sea $r$ un nodo de barra interior con desplazamiento $\mathbf{u}_r$, y
sea $e$ el hexaedro huésped con nodos $\{i\}$ y funciones de forma
$N_i(\xi, \eta, \zeta)$.  El desplazamiento interpolado del hexaedro
en la posición de la barra es:
%
\begin{equation}
  \tilde{\mathbf{u}} = \sum_i N_i\,\mathbf{u}_i.
\end{equation}

Se define un resorte de penalidad con rigidez $\alpha$ que impone:
%
\begin{equation}
  \mathbf{f}_r^{\text{pen}} =
    \alpha\,(\mathbf{u}_r - \tilde{\mathbf{u}}),
  \qquad
  \mathbf{f}_i^{\text{pen}} =
    -\alpha\,N_i\,(\mathbf{u}_r - \tilde{\mathbf{u}}).
\end{equation}

La contribución a la matriz tangente es:
%
\begin{equation}
  \mathbf{K}^{\text{pen}} =
  \begin{bmatrix}
    \alpha\,\mathbf{I}_3 & -\alpha\,N_1\,\mathbf{I}_3 &
      \cdots & -\alpha\,N_n\,\mathbf{I}_3 \\[4pt]
    -\alpha\,N_1\,\mathbf{I}_3 &
      \alpha\,N_1 N_1\,\mathbf{I}_3 & \cdots &
      \alpha\,N_1 N_n\,\mathbf{I}_3 \\
    \vdots & \vdots & \ddots & \vdots \\
    -\alpha\,N_n\,\mathbf{I}_3 &
      \alpha\,N_n N_1\,\mathbf{I}_3 & \cdots &
      \alpha\,N_n N_n\,\mathbf{I}_3
  \end{bmatrix},
  \label{eq:penalty-stiffness}
\end{equation}
%
donde $n$ es el número de nodos del hexaedro (8, 20 o 27) y donde
sólo se incluyen los pares con $|N_i| > 10^{-15}$.

\subsection{Selección de $\alpha$}

En la formulación de penalidad, el error relativo de acoplamiento es
del orden de $k_{\text{elem}} / \alpha$, donde $k_{\text{elem}}$ es
la rigidez diagonal dominante del grado de libertad.  Para los
parámetros del ejemplo:
%
\begin{align}
  k_{\text{truss}} &= E_s\,A_b / L_e
    = 200\,000 \times 3{.}14 \times 10^{-4} / 0{.}0625
    \approx 1\,005\;\text{MPa\,m}, \\[4pt]
  k_{\text{hex}} &\approx E_c \cdot \Delta x \cdot \Delta y
    \approx 30\,000 \times 0{.}25^2
    = 1\,875\;\text{MPa\,m}.
\end{align}

Con $\alpha = 10^6$ (valor por defecto), el error relativo es
$\sim 10^{-3}$, suficiente para la precisión de un modelo de
sub-estructura.

\subsection{Implementación en SNES}

Las contribuciones de penalidad se inyectan directamente en los
callbacks estáticos de PETSc:

\begin{itemize}
  \item \textbf{FormResidual}: Los términos $\mathbf{f}^{\text{pen}}$
        se suman al vector de fuerzas internas en espacio local
        (\texttt{f\_int\_local}) antes del scatter
        \texttt{DMLocalToGlobal}.  Los grados de libertad de nodos
        con Dirichlet se eliminan automáticamente durante el scatter.

  \item \textbf{FormJacobian}: Los términos de
        $\mathbf{K}^{\text{pen}}$ se insertan en la matriz global
        $\mathbf{J}$ mediante \texttt{MatSetValues} con modo
        \texttt{ADD\_VALUES}.  Se utiliza la \emph{global section}
        para obtener los offsets de cada nodo; los nodos restringidos
        (con 0~DOF en la sección global) se omiten.  La opción
        \texttt{MAT\_NEW\_NONZERO\_ALLOCATION\_ERR = PETSC\_FALSE}
        permite que PETSc asigne las nuevas entradas no-cero sin error.
\end{itemize}

Los datos de acoplamiento se pre-calculan una sola vez en
\texttt{compute\_penalty\_couplings()}, invocado después de
\texttt{model\_->setup()} en \texttt{first\_solve()}.  Para cada
nodo de barra interior:

\begin{enumerate}
  \item Se identifica el hexaedro huésped $(i_x, i_y, i_z)$ y las
        coordenadas padres $(\xi, \eta, \zeta)$ desde
        \texttt{RebarNodeEmbedding}.

  \item Se evalúan las funciones de forma 1-D de Lagrange (lineales
        para Hex8, cuadráticas para Hex27) como producto tensorial:
        \[
          N_{p,q,r}(\xi,\eta,\zeta) =
            L_p(\xi) \cdot L_q(\eta) \cdot L_r(\zeta),
        \]
        donde para el caso cuadrático:
        \[
          L_0(t) = \frac{t(t-1)}{2},\quad
          L_1(t) = 1 - t^2,\quad
          L_2(t) = \frac{t(t+1)}{2}.
        \]

  \item Se almacenan los pares $(\text{sieve\_point}_i, N_i)$ en
        un vector \texttt{hex\_weights} dentro de
        \texttt{PenaltyCouplingEntry}.
\end{enumerate}


% ------------------------------------------------------------------
\section{Feedback Fibra-Truss implícito}
\label{sec:emb-feedback}

Una pregunta natural de diseño es cómo retroalimentar la respuesta
del acero del sub-modelo 3D al modelo de fibras de la viga
estructural.  En el esquema de acoplamiento implementado, esta
retroalimentación es \emph{implícita}:

\begin{enumerate}
  \item Las fuerzas seccionales homogeneizadas se calculan mediante
        \texttt{section\_forces\_from\_reactions()}, que suma las
        reacciones internas en \textbf{todos} los nodos de la cara
        $z_{\max}$ (Face~B), incluyendo los nodos de barra que ahora
        están en las listas de condiciones de contorno.

  \item Las fuerzas de los \texttt{TrussElement} contribuyen
        directamente al vector de fuerzas internas del sub-modelo.
        A través del acoplamiento por penalidad, la rigidez del acero
        también afecta la distribución de desplazamientos de los
        hexaedros vecinos.

  \item La tangente seccional $\mathbf{D}_{\text{hom}}$ (6$\times$6)
        se obtiene por perturbación numérica de las componentes
        cinemáticas (deformación axial, curvaturas, cortes,
        torsión).  Cada perturbación modifica las condiciones de
        contorno de las barras de refuerzo en ambas caras, lo cual
        altera el estado de los \texttt{TrussElement} y se refleja
        en las reacciones observadas en Face~B.
\end{enumerate}

Así, el acero del sub-modelo influye tanto en las fuerzas
seccionales como en la rigidez tangente que la viga emplea para el
equilibrio global.  No se requiere una transferencia explícita de
deformaciones del truss a la fibra; el esquema de acoplamiento
estequiométrico FE\textsuperscript{2} (downscale de cinemática,
upscale de fuerzas y tangente) captura automáticamente la
contribución del refuerzo.

\subsection{Extracción de deformaciones axiales del acero}

Para diagnóstico y post-procesamiento, el método
\texttt{extract\_rebar\_strains()} calcula la deformación axial
promedio de cada barra:
%
\begin{equation}
  \overline{\varepsilon}_b =
    \frac{1}{n_z} \sum_{e=1}^{n_z} \varepsilon_{b,e}^{\text{GP}},
\end{equation}
%
donde $\varepsilon_{b,e}^{\text{GP}}$ es la deformación del punto
de Gauss del $e$-ésimo elemento truss de la barra $b$.  Esto permite
comparar la deformación del acero del sub-modelo con la del
correspondiente \emph{fiber} en la sección estructural.


% ------------------------------------------------------------------
\section{Visualización VTK de barras con espesor}
\label{sec:emb-vtk}

La exportación VTK de las barras se extiende con cuatro campos por
celda:

\begin{center}
\begin{tabular}{llp{7cm}}
\toprule
Array & Tipo & Descripción \\
\midrule
\texttt{axial\_stress} & Cell & Tensión axial $\sigma$ [MPa] del
  punto de Gauss del \texttt{TrussElement} \\
\texttt{axial\_strain} & Cell & Deformación axial
  $\varepsilon$ \\
\texttt{bar\_area}      & Cell & Área de la barra $A_b$
  [m$^2$] \\
\texttt{TubeRadius}     & Cell & Radio del tubo
  $\phi_b / 2$ [m] para el filtro \emph{Tube} de ParaView \\
\texttt{displacement}   & Point & Desplazamiento nodal
  $\mathbf{u}$ [m] (3 componentes) \\
\bottomrule
\end{tabular}
\end{center}

En ParaView, al aplicar el filtro \textit{Filters $\to$
Alphabetical $\to$ Tube}, se selecciona \texttt{TubeRadius} como
\emph{Scalars} y se activa \emph{Vary Radius by Absolute Scalar}.
Esto produce tubos cuyo radio coincide con el radio real de la barra,
permitiendo visualizar la armadura con su geometría física correcta
superpuesta a la malla de concreto.


% ------------------------------------------------------------------
\section{Optimizaciones de rendimiento acumuladas}
\label{sec:emb-performance}

A lo largo de las iteraciones previas se implementaron varias
optimizaciones que mejoran significativamente el rendimiento del
análisis multiescala:

\begin{description}
  \item[Cache de nodos de cara:]
    Los vectores \texttt{face\_min\_z\_ids} y
    \texttt{face\_max\_z\_ids} se calculan una sola vez con
    \texttt{cache\_face\_nodes()} y se reutilizan en cada paso
    global (update\_kinematics y compute\_homogenized\_tangent),
    evitando ${\sim}6$ evaluaciones de \texttt{nodes\_on\_face()}
    por paso iterativo.

  \item[Colección diferida de fisuras:]
    La función \texttt{collect\_crack\_data()} se invoca una sola
    vez por paso global (\texttt{end\_of\_step()}) en lugar de en
    cada iteración estancada, eliminando ${\sim}5$ colecciones
    redundantes por paso.

  \item[Filtro de apertura mínima de fisura:]
    Las fisuras con apertura inferior a 0{.}5\,mm/m se excluyen de
    la exportación VTK, reduciendo el número de planos de fisura de
    $\sim$1264 a $\sim$494 (reducción del 61\%).

  \item[Sub-stepping adaptativo verboso:]
    El módulo de sub-stepping adaptativo con bisección registra
    cada sub-paso con tamaño de fracción, número de iteraciones
    SNES y razón de convergencia, permitiendo diagnóstico en tiempo
    real.

  \item[$n_z = 8$ sub-divisiones en $z$:]
    Se aumentó la resolución axial de $n_z = 4$ a 8 para capturar
    mejor la variación de fisuras a lo largo del elemento viga.
\end{description}


% ------------------------------------------------------------------
\section{Resumen de los cambios de código}
\label{sec:emb-code-summary}

\begin{center}
\small
\begin{tabular}{lp{8cm}}
\toprule
Archivo & Cambio principal \\
\midrule
\texttt{PrismaticDomainBuilder.hh} &
  \texttt{RebarBar} con coordenadas físicas $(l_y, l_z)$ y diámetro.
  \texttt{RebarNodeEmbedding} con $(i_x, i_y, i_z, \xi, \eta, \zeta)$.
  Nodos independientes de barra en posiciones exactas. \\[2pt]
\texttt{MultiscaleCoordinator.hh} &
  Almacena embeddings y diámetros.  Agrega nodos de barra en caras
  $z$ a las listas de nodos de contorno. \\[2pt]
\texttt{FieldTransfer.hh} &
  \texttt{SubModelSpec::RebarBar} extendido con diámetro. \\[2pt]
\texttt{NonlinearSubModelEvolver.hh} &
  \texttt{PenaltyCouplingEntry} y
  \texttt{compute\_penalty\_couplings()}.
  Contribuciones de penalidad en \texttt{FormResidual} y
  \texttt{FormJacobian}.  Exportación VTK con \texttt{TubeRadius},
  \texttt{axial\_strain}.  Método
  \texttt{extract\_rebar\_strains()}. \\[2pt]
\texttt{main\_lshaped\_multiscale\_16storey.cpp} &
  Barras con diámetro $\phi = 20$\,mm y posiciones dentro del
  recubrimiento ($c + \phi/2$). \\
\bottomrule
\end{tabular}
\end{center}


% ------------------------------------------------------------------
\section{Conclusiones}
\label{sec:emb-conclusions}

La formulación de acero embebido con acoplamiento por penalidad
resuelve los dos problemas fundamentales del esquema anterior:

\begin{enumerate}
  \item Las barras de refuerzo quedan posicionadas en sus ubicaciones
        físicas reales, respetando el recubrimiento de concreto
        ($c + \phi/2$ desde el borde de la sección).

  \item El acoplamiento Master-Slave por penalidad permite que los
        nodos de las barras sigan la cinemática del continuo
        hexaédrico (perfecta adherencia) sin compartir nodos, lo
        que habilita futuros desarrollos como modelos de
        deslizamiento barra--concreto.
\end{enumerate}

El feedback fibra-truss es implícito: la contribución del acero se
captura a través de las fuerzas seccionales y la tangente
homogeneizada.  La visualización con tubos VTK permite verificar
visualmente la posición y el estado tensional de cada barra dentro
del sub-modelo de concreto.
